%----------
%--- METODOLOGIA

\chapter{METODOLOGIA}
\thispagestyle{myheadings}

Este capítulo apresenta os procedimentos metodológicos adotados no desenvolvimento desta pesquisa.

\section{Classificação da pesquisa}

Esta pesquisa classifica-se como aplicada, pois tem como objetivo gerar 
conhecimentos voltados à solução de um problema prático e específico no contexto da 
SEFAZ-PB. Nesse sentido, foi desenvolvido um novo sistema de ETL para o projeto BDMalhas, 
fundamentado em uma nova arquitetura baseada em tecnologias modernas do ecossistema de Big Data, 
com o objetivo de superar as limitações de desempenho e escalabilidade da solução legada.

Em relação aos objetivos da pesquisa, podemos a caracteriza-la como descritiva e exploratória. 
Seu caráter descritivo se deve pela comparação entre a solução desenvolvida e a solução legada, 
com base em métricas de desempenho, manutenibilidade e utilização de espaço em disco obtidas 
durante os experimentos. Já o caráter exploratório está relacionado à investigação da aplicação 
de uma nova arquitetura para a reestruturação dos processos de ETL na SEFAZ-PB, 
buscando compreender seus benefícios e limitações.

Quanto ao método da pesquisa, adota-se o método hipotético-dedutivo, com abordagem quantitativa. 
Uma vez que parte-se da hipótese de que a reestruturação da arquitetura resultará em melhorias 
significativas de desempenho, escalabilidade e manutenibilidade. Por sua vez, a pesquisa 
caracteriza-se como quantitativa por fundamentar-se na coleta e análise de dados numéricos 
obtidos durante os experimentos, tais como tempos de execução, métricas de manutenibilidade e 
utilização de espaço em disco.

Por fim, a pesquisa é conduzida por meio de um estudo de caso, tomando como objeto de estudo 
o projeto BDMalhas, desenvolvido pelo NUTES em parceria com a SEFAZ-PB. Essa classificação 
decorre do fato de que a pesquisa busca avaliar os efeitos da reestruturação dos processos 
de ETL em um contexto real e específico.

\section{A concretização da solução}

Após o levantamento da arquitetura existente e a identificação de suas limitações, iniciou-se o 
desenvolvimento da nova solução de ETL proposta neste trabalho.

Para isso, definiu-se uma nova arquitetura baseada em tecnologias do ecossistema de Big Data, 
adotando o Apache Airflow como ferramenta responsável pela orquestração dos processos de ETL, o Apache 
Spark como mecanismo de processamento distribuído, um Data Lake em PostgreSQL como fonte dos dados e o 
ClickHouse como banco de dados analítico para a persistência das informações processadas.

Com a arquitetura definida, iniciou-se à implementação dos pipelines de ETL para o projeto BDMalhas. 
Durante essa etapa, buscou-se estruturar os pipelines de forma modular, favorecendo a reutilização 
de componentes e facilitando sua manutenção e evolução.

\section{Procedimentos experimentais e análise}

Após a implementação dos novos pipilines de ETL pra o BDmalhas, foram coletadas métricas como o número de 
nós, a quantidade de dependências entre tarefas e a profundidade dos fluxos, com o objetivo de comparar a 
complexidade estrutural e a manutenibilidade de ambas as arquiteturas.

Essa coleta foi realizada por meio da contagem individual dos elementos presentes nos fluxos correspondentes
a cada um dos quatro documentos fiscais processados pelo BDMalhas, tanto na solução baseada em SSIS quanto 
na arquitetura proposta baseada em Apache Airflow.

Após isso foram realizados experimentos com os quatro tipos de documentos 
fiscais processados pelo BDMalhas: NFe, NFCe, CTe e EFD. Para cada documento, foram executadas 
cargas referentes aos períodos de sete dias, um mês e três meses, compreendidos entre 
janeiro e março de 2020, comparando o desempenho da arquitetura proposta com a solução legada 
baseada em SSIS.

Por fim, foram coletadas métricas referentes ao tempo de execução das cargas, à quantidade de 
registros processados em cada experimento e ao tamanho final das bases analíticas. Os resultados
obtidos foram organizados em tabelas e gráficos, possibilitando a comparação do desempenho da 
arquitetura proposta em relação à solução legada baseada em SSIS.